% Options for packages loaded elsewhere
\PassOptionsToPackage{unicode}{hyperref}
\PassOptionsToPackage{hyphens}{url}
\documentclass[
  12pt,
]{ctexart}
\usepackage{xcolor}
\usepackage{amsmath,amssymb}
\setcounter{secnumdepth}{-\maxdimen} % remove section numbering
\usepackage{iftex}
\ifPDFTeX
  \usepackage[T1]{fontenc}
  \usepackage[utf8]{inputenc}
  \usepackage{textcomp} % provide euro and other symbols
\else % if luatex or xetex
  \usepackage{unicode-math} % this also loads fontspec
  \defaultfontfeatures{Scale=MatchLowercase}
  \defaultfontfeatures[\rmfamily]{Ligatures=TeX,Scale=1}
\fi
\usepackage{lmodern}
\ifPDFTeX\else
  % xetex/luatex font selection
\fi
% Use upquote if available, for straight quotes in verbatim environments
\IfFileExists{upquote.sty}{\usepackage{upquote}}{}
\IfFileExists{microtype.sty}{% use microtype if available
  \usepackage[]{microtype}
  \UseMicrotypeSet[protrusion]{basicmath} % disable protrusion for tt fonts
}{}
\makeatletter
\@ifundefined{KOMAClassName}{% if non-KOMA class
  \IfFileExists{parskip.sty}{%
    \usepackage{parskip}
  }{% else
    \setlength{\parindent}{0pt}
    \setlength{\parskip}{6pt plus 2pt minus 1pt}}
}{% if KOMA class
  \KOMAoptions{parskip=half}}
\makeatother
\setlength{\emergencystretch}{3em} % prevent overfull lines
\providecommand{\tightlist}{%
  \setlength{\itemsep}{0pt}\setlength{\parskip}{0pt}}
\usepackage{bookmark}
\IfFileExists{xurl.sty}{\usepackage{xurl}}{} % add URL line breaks if available
\urlstyle{same}
\hypersetup{
  hidelinks,
  pdfcreator={LaTeX via pandoc}}

\author{SCX}
\date{2026}

\begin{document}

\section{SCX 决策日志}\label{scx-ux51b3ux7b56ux65e5ux5fd7}

\subsection{2026-06-25}\label{section}

\subsubsection{决策 1：第 4 篇独立为 SCX
项目}\label{ux51b3ux7b56-1ux7b2c-4-ux7bc7ux72ecux7acbux4e3a-scx-ux9879ux76ee}

\textbf{决策}：将第 4 篇论文从 EGP 中独立出来，在
\texttt{G:\textbackslash{}Xiaogan\_Supercomputing\_data\textbackslash{}SCX\textbackslash{}}
建立独立项目。

\textbf{Why}： - 第 4 篇本质是数学理论/ML 方法论文，不依赖 DFT 计算 -
独立项目便于专注理论开发、合成实验和 ML benchmark 验证 - 第 1-3 篇继续在
EGP 中迭代（材料势函数主线） - 两线可以并行推进

\subsubsection{决策 2：SCX
不应写成大模型工程论文}\label{ux51b3ux7b56-2scx-ux4e0dux5e94ux5199ux6210ux5927ux6a21ux578bux5de5ux7a0bux8bbaux6587}

\textbf{决策}：第 4 篇定位为''数学框架 + 轻量验证''，不要求大规模 LLM
实验。

\textbf{Why}： - 博士阶段资源有限，无法负担 LLM 标注/训练实验 -
数学框架本身是核心贡献，合成实验足以验证理论 -
大模型数据治理是应用前景，不是实证主战场 - 保留向顶刊/顶会投稿的灵活性

\subsubsection{决策 3：Paper 4 可在 Paper 1
后立即启动}\label{ux51b3ux7b56-3paper-4-ux53efux5728-paper-1-ux540eux7acbux5373ux542fux52a8}

\textbf{决策}：不等 Paper 2/3 完成，即可并行推进 SCX。

\textbf{Why}： - SCX 的核心思想来自 Paper 1 的 gauge fixing 经验 -
合成实验不需要新 DFT 数据 - 只需 Paper 1 有 arXiv preprint
即可作为动机引用 - 可以先写数学框架占坑

\subsubsection{决策 4：核心命名定为
SCX}\label{ux51b3ux7b56-4ux6838ux5fc3ux547dux540dux5b9aux4e3a-scx}

\textbf{决策}：框架命名为 \textbf{SCX: State-Conditioned eXpertise}。

\textbf{Why}： - 短、好记、像算法名 - State-Conditioned =
核心思想（状态条件） - eXpertise = 专家能力/可靠性/专业性 - 有个人印记 -
可以扩展为 SCX reliability、SCX data value、SCX expert routing

\subsection{2026-06-26：确立 5
篇论文谱系}\label{ux786eux7acb-5-ux7bc7ux8bbaux6587ux8c31ux7cfb}

\textbf{决策}：论文序列调整为 EGP Paper 1 → SCX-Theory → SCX-MLIP →
SCX-Sim → SCX-Health

\textbf{理由}：先 concrete（EGP 有完整 DFT 数据）打根据地，再
abstract（SCX-Theory）建理论， 再 apply
back（SCX-MLIP）验证，再跨领域扩展（SCX-Sim, SCX-Health）。

\subsection{2026-06-26：确立通用模块化架构（v4.0
方向）}\label{ux786eux7acbux901aux7528ux6a21ux5757ux5316ux67b6ux6784v4.0-ux65b9ux5411}

\textbf{决策}：SCX v4 进行通用模块化重构，核心是''抽象接口 + Encoder
插件 + YAML 配置''三层分离。

\textbf{来源}：GPT 讨论 (\texttt{与gpt的对话202606261332.md})
中用户指出''框架还是需要更通用一些，避免一直加模块迭代麻烦，应该模块化，不同的行业可以调用相同的代码''。

\textbf{核心设计}：

\begin{enumerate}
\def\labelenumi{\arabic{enumi}.}
\tightlist
\item
  \textbf{SCXStateEncoder / SCXExpert / SCXDataPoint 抽象基类} ---
  所有核心模块只依赖抽象接口
\item
  \textbf{Encoders 插件式} --- 每个领域继承抽象基类，实现自己的 encoder
\item
  \textbf{Config 声明式} --- YAML
  描述完整任务，\texttt{scx\ run\ -\/-config\ domains/mlip.yaml}
\end{enumerate}

\textbf{四个阶段}： - Phase A (2-4周, P0): 核心重构，336 tests 通过 -
Phase B (2-3周, P1): 声明式配置 + CLI - Phase C (4-8周, P1):
新领域扩展（Robot/FEM/LLM） - Phase D (长期, P2): 插件系统 + Benchmark +
社区

\textbf{Why}： - 每加一个领域写一个 adapter，代码重复严重 -
接口不统一（MLIP encoder 和 Vision encoder 完全不同） -
配置散落各处，无法通过 YAML 完整描述一个 SCX 任务 -
商业上无法快速生成''试点→审计报告''的标准流程 -
重构后新增领域的适配成本从''2-3 周''降到''2-3 天''

\textbf{关键文件}： - \texttt{SCX\_规划\_v4\_通用模块化.md} --- 完整规划
- \texttt{SCX\_TODO\_v4\_通用模块化.md} --- 详细待办（替代
\texttt{SCX\_TODO\_2026-06-26.md}）

\subsection{2026-06-27：定理证明里程碑}\label{ux5b9aux7406ux8bc1ux660eux91ccux7a0bux7891}

\subsubsection{决策 1：方案 D --- Theorem 1+2
两阶段发布}\label{ux51b3ux7b56-1ux65b9ux6848-d-theorem-12-ux4e24ux9636ux6bb5ux53d1ux5e03}

\textbf{决策}：Theorem 1 (Noise Detection) 和 Theorem 2 (Weak Feature
Failure) 采用两阶段发布策略：arXiv 先发草稿版 → 经同行审查后发修订版。

\textbf{Why}： - 两个定理的数学根基已经扎实（Chernoff bound、Fano
inequality） - 但需要更多实验验证才能确保 publishing-ready - arXiv
占坑优先于完美

\subsubsection{决策 2：A6 假设添加到 Theorem
1}\label{ux51b3ux7b56-2a6-ux5047ux8bbeux6dfbux52a0ux5230-theorem-1}

\textbf{决策}：为保证 Theorem 1 的方向正确，新增假设
A6：专家在噪声样本上的一致性显著高于在清洁样本上的一致性。

\textbf{Why}： - Lemma 3 的原结构缺少方向性保证 - A6 使 Chernoff bound
的 KL 散度方向判断成立 -
该假设在实践中有直观合理性（多头对同一错标达成一致的概率高于随机）

\subsubsection{决策 3：V(s) 公式
deprecated}\label{ux51b3ux7b56-3vs-ux516cux5f0f-deprecated}

\textbf{决策}：旧数据价值公式 V(s) = rho · r\_bar · L · (1-D) · max SCX
标记为 deprecated，不再推荐使用。

\textbf{Why}： - 该公式本质是启发式（heuristic），缺乏严格的定理保证 -
Theorem 1-3 提供了更well-defined的替代方案 - 保留代码兼容但文档中标记为
deprecated - 未来版本可能完全移除

\subsubsection{决策 4：4→5
篇论文结构拆分}\label{ux51b3ux7b56-445-ux7bc7ux8bbaux6587ux7ed3ux6784ux62c6ux5206}

\textbf{决策}：原 Paper 4 (SCX-Theory) 拆分为 Paper 4 (噪声检测理论) +
Paper 5 (压缩理论)。

\textbf{Why}： - 原 Paper 4 内容过多（3 定理 + 3 命题 + 算法 + 实验） -
噪声检测和压缩理论的读者群/评审群不同 -
两篇论文各自可以更深入地展开数学细节 - 引用链：Paper 4 → Paper
5（压缩依赖噪声检测的结果）

\textbf{各篇论文内容}： - \textbf{Paper 4}: Theorem 1 (噪声检测),
Theorem 2 (弱特征失效), Theorem 3 (不可区分性), Proposition 1'
(全局后悔下界) - \textbf{Paper 5}: Proposition 3' (状态条件加权),
Proposition 4 (压缩保真度), Governance Protocol, low-rank 条件 -
\textbf{Paper 4 → 5}: 依赖关系 --- Paper 5 的压缩定理依赖 Paper 4
的一致性噪声检测

\subsubsection{决策
5：雅洁（yajie）模块命名}\label{ux51b3ux7b56-5ux96c5ux6d01yajieux6a21ux5757ux547dux540d}

\textbf{决策}：数据清理器模块命名为
``yajie.py''（雅洁------优雅的数据清理）。

\textbf{渊源}： - 来自用户女友名字的音译 - ``雅'' = 优雅 (elegant) +
``洁'' = 清洁 (clean) - 功能：定理驱动的数据清理器，利用 Theorem 1-3
进行噪声检测和自动清理

\subsubsection{决策 6：Arrow analogy
移除}\label{ux51b3ux7b56-6arrow-analogy-ux79fbux9664}

\textbf{决策}：从理论文档中完全移除 Arrow's impossibility
analogy，归档到 \texttt{theory/archive/arrow\_analogy\_removed.md}。

\textbf{Why}： - Arrow 定理（社会选择理论）与 SCX 的类比不够精确 -
可能误导读者以为 SCX 试图解决不可能的/impossibility 问题 -
保持理论文档的精确性和专业性

\subsubsection{决策 7：Spring
模块规划（三义）}\label{ux51b3ux7b56-7spring-ux6a21ux5757ux89c4ux5212ux4e09ux4e49}

\textbf{决策}：规划 Spring（Spring = Spring）模块，命名
``三义''（取其''三种正义''之意），与''雅洁''形成对仗。

\textbf{Why}： - ``雅洁''是数据清理（去噪 / denoising） -
``三义''是状态路由（多专家仲裁 / arbitration） - 两个模块共同构成 SCX
的完整数据治理体系

\subsection{2026-06-28：论文 LaTeX
编译环境搭建与修复}\label{ux8bbaux6587-latex-ux7f16ux8bd1ux73afux5883ux642dux5efaux4e0eux4feeux590d}

\subsubsection{决策 1：VS Code LaTeX 插件从 vscode-latex-runner 迁移到
LaTeX
Workshop}\label{ux51b3ux7b56-1vs-code-latex-ux63d2ux4ef6ux4ece-vscode-latex-runner-ux8fc1ux79fbux5230-latex-workshop}

\textbf{决策}：卸载 \texttt{finlay-ab.vscode-latex-runner}，安装
\texttt{james-yu.latex-workshop} v10.16.1。

\textbf{Why}： - \texttt{vscode-latex-runner}
只能手动编译（Ctrl+Alt+B），无法保存即编译 - LaTeX Workshop 是 VS Code
标准 LaTeX 插件（10M+ 安装量），支持 auto-build on save、内置 PDF
预览、SyncTeX - 同时卸载冲突的 \texttt{mathematic.vscode-pdf}（与 LaTeX
Workshop 争抢 PDF 文件打开权）

\subsubsection{决策 2：编译方案选择 pdflatex x 3（非
latexmk）}\label{ux51b3ux7b56-2ux7f16ux8bd1ux65b9ux6848ux9009ux62e9-pdflatex-x-3ux975e-latexmk}

\textbf{决策}：默认编译方案使用三次
\texttt{pdflatex}（直接调用），不使用 \texttt{latexmk} 包装。

\textbf{Why}： - 当前 MiKTeX 缺少 Perl 运行环境，\texttt{latexmk} 不可用
- \texttt{pdflatex} 是 APS/revtex 期刊的标准编译引擎 -
三次编译确保交叉引用和参考文献正确解析 - 禁用 autoClean（同样依赖
latexmk）

\textbf{后续}：如需安装 Perl 恢复 latexmk，MiKTeX 控制台可一键安装。

\subsubsection{决策 3：修复论文代码中的 TeX
非法命令名}\label{ux51b3ux7b56-3ux4feeux590dux8bbaux6587ux4ee3ux7801ux4e2dux7684-tex-ux975eux6cd5ux547dux4ee4ux540d}

\textbf{决策}：将 \texttt{egp\_merging/main.tex} 中
\texttt{\textbackslash{}c0} 全文重命名为 \texttt{\textbackslash{}cb}（25
处），避免与 TeX 内置 \texttt{\textbackslash{}c}（cedilla）命令冲突。

\textbf{Why}： - TeX 控制序列名只能包含字母（A-Z a-z） - \texttt{0}
不是字母，\texttt{\textbackslash{}c0} 被 TeX 解析为
\texttt{\textbackslash{}c} + 字符 \texttt{0} - 这个问题在 Overleaf（TeX
Live）上不报错但本地 MiKTeX 报错------编译器差异

\subsubsection{决策
4：宏双下标问题采用直接展开策略}\label{ux51b3ux7b56-4ux5b8fux53ccux4e0bux6807ux95eeux9898ux91c7ux7528ux76f4ux63a5ux5c55ux5f00ux7b56ux7565}

\textbf{决策}：\texttt{\textbackslash{}cZ\_i} 等 6 处 double subscript
不修改宏定义，改为直接展开为
\texttt{\textbackslash{}mathbf\{c\}\_\{Z\_i\}}。

\textbf{Why}： -
修改宏定义为接受可选参数会引入新的复杂性（\texttt{\textbackslash{}cZ}
本身与 \texttt{\textbackslash{}c} 的命名冲突） - 6
处改动量小，直接展开最安全 - 保留宏在无需物种下标的位置（如
\texttt{\textbackslash{}cZ\^{}\{(m)\}}）继续使用

\subsubsection{相关产出}\label{ux76f8ux5173ux4ea7ux51fa}

\begin{itemize}
\tightlist
\item
  \texttt{.vscode/settings.json}：LaTeX Workshop 完整配置
\item
  \texttt{egp\_merging/main.tex}：4 类 bug 修复，编译通过（17 页）
\item
  \texttt{scx\_method/main.tex}：2 类 bug 修复，编译通过（11 页）
\item
  \texttt{SCX\_发展史与成就.md}：新增第十一节记录
\end{itemize}

\subsubsection{决策 5：确立''作者模式 +
审稿模式''双轨验证方法论}\label{ux51b3ux7b56-5ux786eux7acbux4f5cux8005ux6a21ux5f0f-ux5ba1ux7a3fux6a21ux5f0fux53ccux8f68ux9a8cux8bc1ux65b9ux6cd5ux8bba}

\textbf{决策}：Hermes（6/27）完成证明后，由 Claude Code（6/28）启动
5-agent 独立对抗性审查，不信任 Hermes 的''全修好了''结论。

\textbf{Why}： - Hermes
工作在''证明模式''（作者视角），目标是完成推导、修复已知 bug - Claude
Code 工作在''审稿模式''（reviewer 视角），目标是找到拒稿理由 -
两种模式的认知偏差不同：作者倾向于确认自己是对的，审稿人倾向于假设作者是错的
- 5 个 agent 互不通信、独立出报告，防止群体思维 -
结果验证了方法论：核心定理（Thm 1-3）PASS，但 SE-1 和 Prop 3 在 hostile
review 下暴露了之前未发现的硬伤

\textbf{经验}： - ``全修好了''从作者视角是诚实的------Hermes
确实修了当时已知的所有问题 - 独立对抗性审查发现了 35
个新缺陷------这不是 Hermes 的失误，而是独立审查的价值 -
学术投稿前，两种模式缺一不可：先作者模式完成推导，再审稿模式逐条攻击 -
这个方法应该固化为 SCX 项目的标准流程：任何投稿前都跑一轮 multi-agent
hostile review

\end{document}
